\uuid{WY4J}
\titre{Espaces supplémentaires/Base/Dimension/Polynômes}
\auteur{hanani ; mzari}
\organisation{exo7}
\licence{CC BY-NC-SA 4.0}
\datecreate{2020-01-20}
\niveau{L1}
\module{Algèbre}
\chapitre{Espace vectoriel}
\sousChapitre{Somme directe}
\difficulte{2}
\theme{Espaces vectoriels supplémentaires, Espaces supplémentaires, Base, Dimension, Polynômes}
\exo7qcmid{lille-2.espaces-vectoriels.055}

\contenu{
\begin{question}
\qclassification{106.03}
\qtags{motcle=Espaces supplémentaires/Base/Dimension/Polynômes}
Dans $\Rr_3[X]$, l'espace des polynômes  à coefficients réels  de degré $\le 3$, on considère les deux sous-espaces vectoriels  :
$$E=\{P = a(X-1)^2 +b(X-1)+c\; ;\; a,b,c\in \Rr\}\;\mbox{ et }\; F= \{P= aX^3 +bX^2\; ; \; a,b \in \Rr\}.$$
Quelles sont les assertions vraies ?
\begin{answers}
\bad{$\dim E= 2$}.
\good{$\dim E\cap F = 1$}.
\bad{$E$ et $F$ sont supplémentaires dans $\Rr_3[X]$}.
\good{$E+F=\Rr_3[X]$}.
\end{answers}
\begin{explanations} Une base de $E$ est $\{P_1,P_2,P_3\}$, où $P_1=1, P_2=X-1$, $P_3=(X-1)^2$ donc $\dim E=3$.
Une base de $F$ est $\{Q_1,Q_2\}$, où $Q_1=X^2$ et $Q_2=X^3$, donc $\dim F=2$. En cherchant les relations de dépendance
entre les polynômes $P_1,P_2,P_3, Q_1$ et $Q_2$, on trouve : $P_1+2P_2+P_3= Q_1$. Par conséquent, $E\cap F = \mbox{Vect} \{Q_1\}$, donc  $E$ et $F$ ne sont pas supplémentaires dans $\Rr_3[X]$. Du théorème de la dimension d'une somme, on déduit que $\dim (E+F) = 4$ et comme $E+F$ est un sous-espace de $\Rr_3[X]$ et $\dim \Rr_3[X]=4$, $E+F= \Rr_3[X]$.
\end{explanations}
\end{question}
}
